\chapter{Moral Psychology and AI}
\label{sec:5}
Chapter~\ref{sec:4} asked what a system needs in order to take in a situation at all, and found that the methods most often proposed for giving it values learn those values from somebody who can change their mind. The other half of the job chapter~\ref{sec:3} hands the two of them is here: how a moral sense gets built, in the only case anyone has been able to study, and what of that is available to a machine.

Everything here is an answer to how something comes to hold a standard when nobody is checking. That is what this chapter means by a moral sense that only activates when watched, by a mentor as a control channel, and by the question of whether a system treats its own errors as data or as verdicts. The chapter follows moral development from its social origins, through the friction between fast judgment and slow, out to the institutions and incentives the whole thing runs inside, and down to motivation and play, which are where the difference shows up between a system that has a standard and a system that has had one installed.
